\documentclass[11pt]{article}
\usepackage{amsmath,amssymb,amsthm,hyperref}
\newtheorem{theorem}{Theorem}
\newtheorem{proposition}[theorem]{Proposition}
\begin{document}
\title{Growing Product Domains and Affine Selector Rigidity}
\author{}
\date{November 24, 2026}
\maketitle

\paragraph{Status.}
The selector rigidity theorem below is kernel checked in
\texttt{Results/LargeDomainSelectors.lean}. The budget analysis is a written
audit, not a new upper bound for the original worst-case head complexity.
The original model, admissibility conditions, and complexity definitions
are unchanged. The analysis compares the checked affine-plane compiler with
the independent-cell argument in \texttt{exponent-mechanisms.tex} at
\texttt{origin/ar/0f08fe6fb98a}.

\section{The size test before constructing a decoder}

Let an independently programmable local domain have $M$ inputs in an
ambient Boolean cube of dimension $n$. The referenced independent-cell
argument gives the necessary condition
\[
 H\leq \frac{CM}{nG(n)},\quad G(n)\longrightarrow\infty
 \quad\Longrightarrow\quad
 \liminf\frac{\log_2 M}{G(n)}\geq\frac1{2C}.
\]
This uses Warren's analytic sign-pattern theorem. It is a written necessary
condition for independent universal cells, not a new unconditional Lean
lower bound. In particular, a fixed-power gain $G(n)=n^\delta$ requires
$M\geq\exp(\Omega(n^\delta))$. An unbounded improvement $g(n)$ beyond the
existing logarithmic saving requires $G(n)=\log n\,g(n)$, hence
$M\geq n^{\Omega(g(n))}$. Every fixed-degree polynomial-size local domain
fails the latter test.

Split $n=km$ coordinates into $k$ blocks. The product of $k$ one-hot sets
has $M=m^k$ points and therefore passes the size test only if
\[
 k\log m=\Omega(n^\delta)
 \quad\hbox{or, respectively,}\quad
 k\log m=\Omega(\log n\,g(n)).
\]
Thus growing products are legitimate candidate domains, whereas every
fixed number of one-hot blocks is ruled out at the requested rates.

\section{Two charges that the larger cardinality does not remove}

Compiling $m^{k-2}$ independent two-dimensional slices with the inherited
affine-plane compiler costs
\[
 O\left(m^{k-2}\frac m{\log m}\right)
 =O\left(\frac{M}{m\log m}\right).
\]
Its inputs-per-head scale is $m\log m=(n/k)\log(n/k)$. The enlarged
domain therefore gives no improvement by this compilation method. This
calculation does not rule out a joint decoder for the full tensor.

There is also a separate cover cost. The inherited formal theorem
\texttt{exists\_openDownwardCover} bounds one block's cover by
$22\cdot2^m/m$ centers. Its direct $k$-fold product has the certified bound
\[
 22^k\frac{2^n}{m^k}.
\]
Even granting an ideal local decoder of cost
$C_0M/(n\log M)$, multiplication of these bounds gives only
\[
 C_0\,22^k\frac{2^n}{n k\log m}.
\]
Relative to $2^n/(n\log n)$, the putative saving factor certified by this
calculation is $k\log m/(22^k\log n)\leq k/22^k$, which cannot diverge.
This is a failure of the inherited cover estimate as a growing-product
proof, not a lower bound on the optimal cover. Axis repair and gluing on
all lower-dimensional boundary faces would be additional charges. No
general $k$-block gluing theorem is asserted here.

\section{The first failed lemma in direct affine-plane enlargement}

For a field $F$, put
$\ell_a(u,v)=v-au$. The indicator of the nonvertical line with slope
$a$ and intercept $b$ is $\mathbf1_{\{\ell_a(u,v)=b\}}$.
For a memory-address map $g:A\times B\to F^2$, call its line selectors
additive if, for every $a,b\in F$, there are real functions $U_{a,b}$,
$V_{a,b}$ and a real constant $c_{a,b}$ such that
\[
 \mathbf1_{\{\ell_a(g(x,y))=b\}}
 =c_{a,b}+U_{a,b}(x)+V_{a,b}(y).
 \tag{1}
\]
These are exactly the values an affine numerator can take on a product
of two one-hot sets. They are not the values of a general rational score.

\begin{theorem}[Three-direction rigidity]
Suppose $A,B$ are nonempty and $F$ contains three distinct slopes.
If (1) holds, then $g$ depends on only one coordinate: either
$g(x,y)=u(x)$ for some $u:A\to F^2$, or $g(x,y)=v(y)$ for some
$v:B\to F^2$. In particular, if both factors have at least two elements,
an injective address assignment cannot satisfy (1).
\end{theorem}
\begin{proof}
Every additive function has zero mixed rectangle difference. Fix a
rectangle and abbreviate its corners by $p=g(x,y)$, $q=g(x,y')$,
$r=g(x',y)$, and $s=g(x',y')$. Apply the rectangle identity to the line
of slope $a$ through $p$. Its left side contains the contribution $1$
from $p$, so at least one of $q,r$ lies on this line. Therefore
\[
 \ell_a(p)=\ell_a(q)\quad\hbox{or}\quad\ell_a(p)=\ell_a(r).
 \tag{2}
\]
Two distinct slopes determine a point: subtracting the two intercept
equalities gives $(a-b)(p_1-q_1)=0$, then $p_1=q_1$ and $p_2=q_2$.
Among three instances of (2), two choose the same neighboring corner.
Consequently every rectangle satisfies $p=q$ or $p=r$.

If every row is constant there is nothing to prove. Otherwise choose
$x_0,y_0,y_1$ with $g(x_0,y_0)\ne g(x_0,y_1)$. For any $y$, its value
$g(x_0,y)$ differs from at least one of these two values. Apply the
corner property with that column and any row $x$. It forces
$g(x,y)=g(x_0,y)$. Thus every column is constant across rows.
Injectivity is then impossible when both factors have two elements.
\end{proof}

The three-direction hypothesis is sharp for this statement. Over
$F=\mathbb F_2$, take $g(x,y)=(x+y,y)$ on $\mathbb F_2^2$. The map is
injective, but its slope-zero intercept is $y$ and its slope-one intercept
is $x$. Every nonvertical line indicator depends on just one coordinate.

The current compiler uses weighted selectors
$\eta_b(a)\mathbf1_{\{\ell_a(g)=c\}}$. For a fixed bit query, three
nonzero weights on distinct slopes suffice for the same proof: divide
the corresponding additive formulas by those weights. This observation
does not assume that all fixed Cauchy residue choices are nonzero.
It also does not cover architectures that change the bit query jointly
with the product coordinates, use nonlinear selector scores, cancel
selector errors globally, or ask only for the final strict signs.

\section{Checked declarations and remaining opportunity}

The namespace is \texttt{HeadComplexity.LargeDomain}. The main declarations
are \texttt{additive\_line\_selectors\_depend\_on\_one\_axis} and
\texttt{not\_additive\_line\_selectors\_of\_injective}. Their assumptions
are explicit, and their proofs use no added axiom or \texttt{sorry}.
The supporting theorem \texttt{depends\_on\_one\_axis\_of\_corner\_property}
is purely set theoretic.

The obstruction identifies the first failed lemma in the direct proposal
to reuse the old affine line-rank factors on a larger Cartesian row
domain. It does not prohibit a large-domain rational decoder with
nonlinear intermediate selectors. The remaining task is to construct
such a decoder at cost $O(M/(n\log M))$ or better, with a cover and
gluing density small enough to preserve the larger logarithm. Interaction
across many original cells is another possible escape because their
individual selectors need not then be represented exactly.
\end{document}
